\subsection{自由模的外代数}

定理1. 设M是一个具有基的A-模 \((e_{\lambda})_{\lambda \in L}\) 。设 L 被赋予全序集的结构（集合论，III，§2，第 3 号，定理 1），并且对于 L 的每个有限子集 J，我们记

\[
e _ {J} = e _ {\lambda_ {1}} \wedge e _ {\lambda_ {2}} \wedge \dots \wedge e _ {\lambda_ {n}}\tag{18}
\]

（其中 \((\lambda_k)_{1 \leqslant k \leqslant n}\) 是 \(J\) 的元素按递增顺序排列的序列（集合论，III，§5，第 3 号，命题 6）；我们记 \(e_\varphi = 1\) ，即 A 的单位元。则诸 \(e_J\) （其中 \(J\) 遍历有限子集组成的集合 \(\mathfrak{F}(L)\) （ \(L\) 的））构成外代数 \(\bigwedge(M)\) 的一组基。

由于 \(e_{\lambda}\) 生成 A-模 M， \(\bigwedge(\mathbf{M})\) 的每个元素都是有限个 \(e_{\lambda}\) 之乘积的线性组合，从而（考虑到第 3 号命题 5）是有限个元素 \(e_{J}\) （ \(J \in \mathfrak{F}(L)\) ）的线性组合。剩下要证明 \(e_{J}\) 在 A 上线性无关。否则，这些元素之间将存在一个系数不全为零的线性关系；此关系中诸 \(e_{J}\) （其系数 \(\neq 0\) 者）所对应的子集 J 的并是 L 的一个有限子集 K（因为系数 \(\neq 0\) 的只有有限个）。设 N 为 M 的由诸 \(e_{\lambda}\) （满足 \(\lambda \in K\) 者）生成的子模；N 是 M 的一个直因子，因此（第 2 号） \(\bigwedge(\mathbf{N})\) 等同于 \(\bigwedge(\mathbf{M})\) 的一个子代数，并且，如果我们证明诸 \(e_{J}\) （ \(J \subset K\) ）构成 \(\bigwedge(\mathbf{N})\) 的一组基，我们将得到所期望的矛盾。

因此，问题归结为在 M 的基为有限时证明定理 1；于是我们可以假设 \(L = [1, m] \subset N\) 。对于每个 \(i \in L\) ，设 \(M_i\) 为 M 的自由子模 \(Ae_i\) ；M 是诸 \(M_i\) 的直和，而 \(\bigwedge(M_i)\) 是 \(\bigwedge^0(M_i) = A\) 与 \(\bigwedge^1(M_i) = M_i\) （第 3 号，命题 6）的直和。令 \(\bigwedge(M)\) 与诸 \(\bigwedge(M_i)\) 的张量积这一 A-模（第 7 号，命题 10）典范地等同；后者以诸基 \((1, e_i)\) （诸 \(\bigwedge(M_i)\) 的基）的张量积为基（II，§ 3，第 7 号，命题 7 的推论 2）；因此我们得到所有元素

\[
u _ {1} \otimes u _ {2} \otimes \dots \otimes u _ {m}
\]

其中 \(u_{i}=1\) 或 \(u_{i}=e_{i}\) ；若 J 是使得 \(u_{i}=e_{i}\) 的指标 i 组成的集合，则 \(u_{1}\otimes u_{2}\otimes\cdots\otimes u_{m}\) 与 \(e_{J}\) 等同，这就完成了证明。

推论 1. 假设 \(\mathbf{L} = [1, m]\) ；则基 \((e_{\mathrm{J}})_{\mathrm{J} \in \mathfrak{P}(\mathrm{L})}\) （ \(\bigwedge (\mathbf{M})\) 的）有 \(2^{m}\) 个元素。对于 \(p > m\) , \(\bigwedge^{p}(\mathbf{M}) = \{0\}\) ； \(\bigwedge^{m}(\mathbf{M})\) 有由单个元素 \(e_{\mathrm{L}}\) 组成的基；对于 \(0 \leqslant p \leqslant m\) ，基 \((e_{\mathrm{J}})\) （ \(\bigwedge^{p}(\mathbf{M})\) 的、由诸 \(e_{\mathrm{J}}\) （满足 \(\operatorname{Card}(\mathbf{J}) = p\) 者）组成的）的元素个数是

\[
\binom {m} {p} = \frac {m !}{p ! (m - p) !}.
\]

这可由《集合论》III，§3，第5号，命题12以及《集合论》III，§5，第8号，命题11的推论1得出。

我们回到定理 1 中集合 L 为任意的情形，并明确给出基 \((e_{J})\) 的乘法表（§1，第 7 号）。给定两个有限子集 \(J\) , \(K\) （全序集 \(L\) 的），我们记

\[
\left\{ \begin{array}{l l} \rho_ {J, K} = 0 & \text { if } J \cap K \neq \varnothing \\ \rho_ {J, K} = (- 1) ^ {\nu} & \text { if } J \cap K = \varnothing \end{array} \right.\tag{19}
\]

其中 v 表示有序对 \((\lambda, \mu) \in J \times K\) 中满足 \(\lambda > \mu\) 者的个数。于是，第 2 号命题 4 的推论 1 立即蕴含关系

\[
e _ {J} \wedge e _ {K} = \rho_ {J, K} e _ {J \cup K}.\tag{20}
\]

注意公式

\[
\rho_ {J, K} \rho_ {K, J} = (- 1) ^ {j k}\tag{21}
\]

当 \(\mathbf{J} \cap \mathbf{K} = \varnothing, j = \operatorname{Card}(\mathbf{J}), k = \operatorname{Card}(\mathbf{K})\) (第3号，命题5之推论2。)

推论2. 若M是投射A-模， \(\bigwedge (\mathbf{M})\) 是投射 A-模。

证明与第5节第5号定理1的推论相同，只需将T替换为 \(\Lambda\) .

推论3. 设 M 为投射 A-模，且 \((x_{i})_{1\leqslant i\leqslant n}\) 为 M 的元素组成的有限序列。要在 M 上存在一个交错 \(n\) -线性形式 \(f\) ，使得

\[
f (x _ {1}, x _ {2}, \dots , x _ {n}) \neq 0,
\]

其充要条件是 \(x_{1} \wedge x_{2} \wedge \cdots \wedge x_{n} \neq 0\) .

我们已经知道（在对 M 不作任何假设的情况下）该条件是必要的（第 4 号，命题 7）。现在假设 M 是投射的，并且

\[
x _ {1} \wedge x _ {2} \wedge \dots \wedge x _ {n} \neq 0.
\]

那么 \(\bigwedge^n (\mathbf{M})\) 是投射的（推论 2），因此典范映射

\[
\bigwedge^ {n} (\mathbf {M}) \rightarrow (\bigwedge^ {n} (\mathbf {M})) ^ {* *}
\]

是单射的（II，§2，第7号，命题13的推论4）；我们由此得出结论，存在一个线性形式 \(g: \bigwedge^{n}(\mathbf{M}) \to \mathbf{A}\) ，使得 \(g(x_{1} \wedge x_{2} \wedge \dots \wedge x_{n}) \neq 0\) 。若 \(f\) 是交错 \(n\) -线性形式中与 \(g\) 相对应的那一个（第 4 号，命题 7），则 \(f(x_{1}, \ldots, x_{n}) \neq 0\) .

